\section{Architektura systemu}

\subsection{Widok komponentów}
System składa się z czterech głównych części uruchamianych przez Docker Compose:
\begin{itemize}
  \item \textbf{Frontend} (Vite + React + TypeScript) -- zbiera telemetrię użytkownika.
  \item \textbf{Backend} (FastAPI) -- kolekcja eventów, predykcja, etykietowanie, API.
  \item \textbf{PostgreSQL} -- trwałe składowanie sesji, eventów, predykcji i akcji.
  \item \textbf{Bot Runner} (Playwright) -- generowanie powtarzalnych scenariuszy botowych.
\end{itemize}

\begin{figure}[H]
  \centering
  \resizebox{0.98\textwidth}{!}{\begin{tikzpicture}[
  x=1cm, y=1cm,
  comp/.style={
    rectangle,
    rounded corners,
    draw=black,
    very thick,
    align=center,
    text width=3.6cm,
    minimum height=1.3cm,
    fill=blue!6
  },
  store/.style={
    cylinder,
    shape border rotate=90,
    draw=black,
    very thick,
    align=center,
    text width=4.6cm,
    minimum height=1.8cm,
    aspect=0.27,
    fill=green!8
  },
  link/.style={-Latex, thick}
]

\node[comp, fill=orange!14] (user) at (0.0, 0.0) {Użytkownik\\(przeglądarka)};
\node[comp] (frontend) at (5.1, 0.0) {Frontend\\Vite + React};
\node[comp] (backend) at (10.6, 0.0) {Backend API\\FastAPI};
\node[store] (db) at (10.6, -3.1) {PostgreSQL\\sessions, events, predictions};
\node[comp, fill=purple!8] (bot) at (5.1, -2.4) {Bot Runner\\Playwright};
\node[comp, fill=teal!8] (ml) at (10.6, 3.2) {Feature Builder +\\Decision Tree + Policy};

\draw[link] (user) -- node[midway, above, fill=white, inner sep=1pt]{interakcje UI} (frontend);
\draw[link] (frontend) -- node[midway, above, fill=white, inner sep=1pt]{batch telemetry} (backend);
\draw[link] (backend) -- node[midway, right, fill=white, inner sep=1pt]{zapis/odczyt} (db);
\draw[link] (backend) -- node[midway, right, fill=white, inner sep=1pt]{okno cech i predykcja} (ml);
\draw[link, dashed] (bot) -- node[midway, left, fill=white, inner sep=1pt]{scenariusze botowe} (frontend);
\draw[link, dashed] (bot.east) to[bend left=14] node[midway, above, fill=white, inner sep=1pt]{etykietowanie botów} (backend.south west);

\end{tikzpicture}
}
  \caption{Architektura logiczna MVP}
\end{figure}

\subsection{Przepływ danych end-to-end}
\begin{enumerate}
  \item Frontend inicjuje sesję przez \texttt{POST /api/v1/sessions}.
  \item Kolektor telemetrii zbiera eventy i wysyła je batchami do \texttt{/api/v1/events/batch}.
  \item Backend zapisuje surowe dane do tabel \texttt{sessions} i \texttt{raw\_events}.
  \item Dla żądania predykcji backend buduje cechy z okna eventów.
  \item Model zwraca \texttt{predicted\_label}, prawdopodobieństwo bota i \texttt{risk\_score}.
  \item Policy engine mapuje \texttt{risk\_score} na akcję enforcement.
  \item Wynik i akcja są zapisywane (\texttt{predictions}, \texttt{enforcement\_actions}).
\end{enumerate}

\subsection{Diagram przypadków użycia}
\begin{figure}[H]
  \centering
  \resizebox{0.95\textwidth}{!}{\begin{tikzpicture}[
  x=1cm, y=1cm,
  actor/.style={draw, circle, minimum size=0.7cm, fill=gray!10},
  uc/.style={
    ellipse,
    draw=black,
    thick,
    align=center,
    minimum height=1.2cm,
    fill=blue!6,
    text width=6.4cm
  },
  line/.style={-Latex, thick}
]

\node[actor] (human) at (0.0,7.4) {};
\node[align=center] at (0.0,6.6) {Człowiek};

\node[actor] (bot) at (0.0,4.2) {};
\node[align=center] at (0.0,3.4) {Bot runner};

\node[actor] (operator) at (0.0,0.9) {};
\node[align=center] at (0.0,0.1) {Analityk};

\node[uc] (u1) at (8.2,7.4) {Generowanie\\telemetrii sesji};
\node[uc] (u2) at (8.2,5.8) {Wysłanie batcha eventów\\do Collector API};
\node[uc] (u3) at (8.2,4.2) {Predykcja human vs bot\\+ risk score};
\node[uc, text width=8.0cm] (u4) at (8.2,2.3) {Decyzja polityki\\allow/observe/throttle/challenge/block};
\node[uc] (u5) at (8.2,0.5) {Etykietowanie sesji\\do treningu};

\draw[line] (human) -- (u1);
\draw[line] (human) -- (u2);
\draw[line] (human) -- (u3);
\draw[line] (bot) -- (u1);
\draw[line] (bot) -- (u2);
\draw[line] (bot) -- (u3);
\draw[line] (operator) -- (u5);
\draw[line] (operator) -- (u4);

\draw[line, dashed] (u1.south) -- (u2.north);
\draw[line, dashed] (u2.south) -- (u3.north);
\draw[line, dashed] (u3.south) -- (u4.north);

\end{tikzpicture}
}
  \caption{Główne przypadki użycia systemu}
\end{figure}

\subsection{Model danych (ERD)}
\begin{figure}[H]
  \centering
  \resizebox{0.98\textwidth}{!}{\begin{tikzpicture}[
  x=1cm, y=1cm,
  tbl/.style={
    rectangle,
    draw=black,
    thick,
    align=left,
    font=\footnotesize,
    fill=green!6,
    text width=4.9cm,
    inner sep=5pt
  },
  rel/.style={-Latex, thick}
]

\node[tbl] (sessions) at (0.0, 4.8) {
\textbf{sessions}\\
PK: id\\
client\_fingerprint (JSON)\\
environment (JSON)\\
source, true\_label\\
started\_at\\
last\_event\_at
};

\node[tbl] (events) at (0.0, 1.0) {
\textbf{raw\_events}\\
PK: id\\
FK: session\_id $\rightarrow$ sessions.id\\
sequence\_no\\
event\_type, ts\_ms\\
x, y, scroll\_x, scroll\_y\\
target\_id, pointer\_type\\
payload
};

\node[tbl] (pred) at (7.0, 4.8) {
\textbf{predictions}\\
PK: id\\
FK: session\_id $\rightarrow$ sessions.id\\
FK: model\_version\_id $\rightarrow$ model\_versions.id\\
predicted\_label\\
probability\_bot\\
risk\_score, confidence\\
created\_at
};

\node[tbl] (modelv) at (7.0, 1.0) {
\textbf{model\_versions}\\
PK: id\\
name, version, algorithm\\
metrics (JSON)\\
artifact\_path\\
metadata (JSON)\\
created\_at
};

\node[tbl] (enf) at (7.0, -2.7) {
\textbf{enforcement\_actions}\\
PK: id\\
FK: session\_id $\rightarrow$ sessions.id\\
FK: prediction\_id $\rightarrow$ predictions.id\\
action, reason\\
created\_at
};

\draw[rel] (events.north) -- node[midway, left, fill=white, inner sep=1pt]{N:1} (sessions.south);
\draw[rel] (pred.west) -- node[midway, above, fill=white, inner sep=1pt]{N:1} (sessions.east);
\draw[rel] (pred.south) -- node[midway, right, fill=white, inner sep=1pt]{N:1} (modelv.north);
\draw[rel] (enf.north) -- node[midway, right, fill=white, inner sep=1pt]{N:1} (pred.south);
\draw[rel] (enf.west) to[bend right=24] node[pos=0.42, left, fill=white, inner sep=1pt]{N:1} (sessions.south east);

\end{tikzpicture}
}
  \caption{Relacyjny model danych PostgreSQL}
\end{figure}
